% !TeX program = xelatex
% 独立可编译源文件。正文起讫位置及需核实事项已用注释标出。
% 编译：xelatex -interaction=nonstopmode 南航开放课题_报告正文.tex（运行两次）
% 项目范围：既定连接构型下，柔性支撑机器人耦合动力学、闭环辨识及运动振动控制。
% 申请金额暂按30万元起草；量化指标均为拟定目标，不是已有实验结果。
\documentclass[UTF8,a4paper,zihao=-4,fontset=fandol]{ctexart}
\usepackage[left=2.6cm,right=2.6cm,top=2.5cm,bottom=2.5cm]{geometry}
\usepackage{amsmath,amssymb,bm,booktabs,longtable,array,tabularx}
\usepackage{enumitem,setspace,xcolor,caption,needspace}
\usepackage{tikz}
\usetikzlibrary{arrows.meta,positioning,calc}
\usepackage[unicode,hidelinks]{hyperref}
\usepackage{xurl}
\hypersetup{pdftitle={柔性支撑空间机器人耦合动力学与运动振动协同控制},pdfsubject={南航重点实验室开放课题申请书报告正文}}
\setstretch{1.32}
\setlength{\parindent}{2em}
\setlength{\parskip}{0.15em}
\setlength{\emergencystretch}{2em}
\clubpenalty=10000
\widowpenalty=10000
\displaywidowpenalty=10000
\raggedbottom
\setlist[enumerate]{leftmargin=2.2em,itemsep=0.3em,topsep=0.3em}
\captionsetup{font=small,labelfont=bf,labelsep=quad,skip=8pt}
\ctexset{
 section={format=\zihao{4}\bfseries,name={,、},number=\chinese{section},aftername={},beforeskip=1.2em,afterskip=0.7em},
 subsection={format=\zihao{-4}\bfseries,number=\arabic{subsection},aftername={.\ },beforeskip=1em,afterskip=0.5em},
 subsubsection={format=\zihao{-4}\bfseries,number=\arabic{subsection}.\arabic{subsubsection},aftername={\quad},beforeskip=0.8em,afterskip=0.3em}
}
\newcommand{\topic}[1]{\par\addvspace{0.5em}\noindent\textbf{#1}\par\nobreak}
\newcommand{\needinfo}[1]{\par\noindent{\color{red!65!black}\textbf{〔待补充〕}#1}\par}
\newcommand{\doi}[1]{\href{https://doi.org/#1}{\nolinkurl{#1}}}
\newcolumntype{Y}{>{\raggedright\arraybackslash}X}
\newcolumntype{P}[1]{>{\raggedright\arraybackslash}p{#1}}
\renewcommand{\arraystretch}{1.25}
\pagestyle{plain}

\begin{document}

% 搬入Word模板时，以下题目填入“项目名称”；从“一”开始搬入报告正文。
\begin{center}
 {\zihao{3}\bfseries 柔性支撑空间机器人耦合动力学\par
 与运动振动协同控制\par}
 \vspace{0.55em}
 {\zihao{-4}航空航天结构力学及控制全国重点实验室开放课题\quad 重点课题}
\end{center}
\vspace{0.35em}

% ==================== 报告正文开始 ====================
\section{立项依据与研究内容}

\subsection{项目的立项依据}

\subsubsection{研究意义}

大型空间天线、空间太阳能电站等装备向大尺度、轻量化发展，使在轨组装成为突破单次发射尺寸约束的重要技术途径。空间机器人依附于已建结构完成模块搬运与安装时，结构承担作业支撑和载荷传递功能，其柔性变形直接参与机器人运动，需要从结构动力学与控制的结合处解决作业精度问题\cite{hu2025,wu2024}。

本项目关注在轨组装中一个具有代表性的作业情形：机器人安装在轻质柔性支撑上，携带模块沿给定轨迹运动并在安装位置附近停驻。在这一过程中，关节加减速产生的反作用载荷传入支撑结构，引起低频振动；支撑点的位移和转动又通过机械臂运动链传递至末端，改变实际运动轨迹。当支撑刚度、阻尼或末端载荷与控制模型存在偏差时，结构响应的幅值与相位预测失准，原有控制参数难以同时保证末端精度和振动衰减。运动结束后的残余振动还将延长下一步操作的等待时间，影响连续作业效率。

这一问题的难点在于，柔性支撑的运动不仅是需要补偿的外部扰动，也是机器人自身动作激发的系统响应。单纯依据末端误差增加控制增益，可能通过反作用载荷进一步激励结构；仅降低运动速度，又难以满足作业效率要求。需要回答：哪些结构模态及耦合参数决定末端误差，哪些参数能够从受控运动中可靠辨识，以及模型在线修正后如何维持运动与振动控制的协调关系。这些问题具有明确的刚柔耦合力学内涵，也是复杂结构闭环辨识与自适应控制的具体应用载体。

本项目拟研究柔性支撑机器人面向控制的动力学表征和参数辨识，发展兼顾轨迹跟踪与结构振动的自适应控制方法，为在轨组装作业参数选择、振动约束设置及控制器设计提供依据。

\subsubsection{国内外研究现状及发展动态}

\topic{（1）空间机器人与柔性结构的耦合动力学建模}

空间机器人多体建模和动量耦合分析已能描述机械臂对承载平台的作用\cite{papadopoulos2021}。对于柔性结构，Sonneville等研究了模态缩减方法，为保持运动学一致性和提高计算效率提供了途径\cite{sonneville2021}。国内研究进一步将模块连接关系引入动力学模型，分析了结构增长对主动控制的影响\cite{wang2020}。

在多柔性杆空间机械臂研究中，Ni等分析了动态边界条件与模态选取的关系，其算例表明，适当的定常模态近似也可能获得满足要求的精度\cite{ni2024}。这说明面向控制的模型选择需要结合具体频段和工况判断，不能仅以模型复杂度衡量有效性。对柔性支撑机器人而言，仍需将结构动力学模型与实际可测量、可调控的变量联系起来，明确支撑响应如何传递至末端，以及模型截断误差如何影响控制效果。

\topic{（2）运动跟踪与结构振动控制}

利用机械臂运动调节柔性基座响应是该领域的重要研究思路。Nenchev等提出反作用零空间控制，将末端运动与柔性基座振动抑制联系起来，说明机器人驱动可同时承担操作和结构调控功能\cite{nenchev1999}。近年来，研究进一步关注末端反馈、柔性状态估计及参数不确定性。Patel和Damaren利用估计的弹性坐标实施末端与振动控制，并通过柔性机械臂数值算例验证了方法\cite{patel2024}；Zhang等研究了参数不确定条件下的自适应鲁棒控制，为运动与振动联合调节提供了参考\cite{zhang2025}。

现有方法为本项目奠定了基础，但其适用条件需要与柔性支撑作业场景对应。机械臂自由度、任务约束和传感配置决定了可用于抑振的控制能力；当末端运动自由度基本占用全部驱动能力时，需要在允许的跟踪误差范围内分配控制作用。结构模态和载荷参数发生变化后，固定补偿模型及固定振动反馈参数的适用范围也需要重新评估。因此，有必要结合参数辨识研究模型变化与控制性能之间的关系。

\topic{（3）闭环动力学辨识与模型修正}

从操作数据中辨识结构与载荷参数，是减轻精确先验模型依赖的重要途径。Ni等采用递归子空间方法辨识柔性空间机械臂的状态空间模型和末端载荷质量，展示了利用输入输出数据修正动力学模型的可行性\cite{ni2019}。闭环辨识理论指出，反馈作用使输入与测量噪声相关，辨识方法需要同时考虑动力学模型、噪声模型及外部激励条件\cite{forssell1999}。将一般离线拟合直接用于闭环运行，可能得到预测误差较小但控制适用性不足的模型。

学习增强控制为处理未被名义模型充分描述的动力学提供了补充手段，相关研究也强调对学习误差、状态约束及控制稳定性进行联合分析\cite{brunke2022}。柔性支撑系统的困难在于，参数偏差、结构振动和测量误差可能在同一频段影响输出。若缺少可辨识性分析和物理约束，在线更新可能将噪声误认为参数变化，或者将真实振动吸收进补偿项。因而需要优先辨识对控制性能敏感的少量参数，并限定数据补偿的作用范围。

\topic{（4）地面验证与本项目的研究切入点}

团队已开展空间机器人模块组装及对接力控制实验\cite{shi2024}，并研究了空间多柔体系统缩比等效建模方法\cite{qi2025}。地面验证需要区分重力、支撑装置附加约束与目标柔性效应，明确实验结果的适用范围。

综合上述进展，本项目将研究重点落在\textbf{柔性支撑响应的可辨识表征及其对运动振动协同控制的作用机制}。在既定连接构型下，以少量主导模态与敏感参数描述作业频段内的耦合特性；通过闭环数据修正模型，并将辨识误差和补偿作用纳入控制分析；最后在相同任务轨迹、执行时间及驱动约束下进行对比实验，检验精度改善是否伴随结构振动的实际降低。

\subsubsection{与申请指南的关系}

本项目主要对应指南2.2“轻质柔性结构动力学设计与分析”，围绕超大尺度柔性空间结构在轨组装场景，研究机器人参与作业时的耦合动力学响应及振动抑制问题；同时对应指南2.5“闭环系统模型辨识、分析与控制”，研究强耦合、多振动模态和参数不确定条件下的闭环辨识与自适应控制。研究对象限定为具有代表性的柔性支撑单元及其承载机器人，采用若干既定支撑状态表征不同作业条件，每次运动试验保持连接关系不变，重点研究模块搬运和接触前停驻阶段的运动及振动响应。该范围与两年开放课题的实施条件相适应。

% 参考文献依模板置于立项依据之后。均核验至期刊、出版商或机构页面。
\begingroup
\small
\setstretch{1.12}
\renewcommand{\refname}{主要参考文献}
\begin{thebibliography}{99}
\setlength{\itemsep}{0.3em}
\bibitem{hu2025} 胡海岩, 田强, 文浩, 罗凯, 马小飞. 极大空间结构在轨组装的动力学与控制[J]. 力学进展, 2025, 55(1): 1--29. DOI: \doi{10.6052/1000-0992-24-044}.
\bibitem{wu2024} 吴志刚, 蒋建平, 邬树楠, 李庆军, 王兴, 谭述君, 邓子辰. 航天结构空间组装动力学与控制研究进展[J]. 力学进展, 2024, 54(2): 344--390. DOI: \doi{10.6052/1000-0992-23-041}.
\bibitem{papadopoulos2021} Papadopoulos E, Aghili F, Ma O, Lampariello R. Robotic manipulation and capture in space: A survey[J]. Frontiers in Robotics and AI, 2021, 8: 686723. DOI: \doi{10.3389/frobt.2021.686723}.
\bibitem{sonneville2021} Sonneville V, Scapolan M, Shan M, Bauchau O A. Modal reduction procedures for flexible multibody dynamics[J]. Multibody System Dynamics, 2021, 51(4): 377--418. DOI: \doi{10.1007/s11044-020-09770-w}.
\bibitem{wang2020} 王恩美, 邬树楠, 吴志刚. 在轨组装空间结构面向主动控制的动力学建模[J]. 力学学报, 2020, 52(3): 805--816. DOI: \doi{10.6052/0459-1879-19-375}.
\bibitem{ni2024} Ni S, Chen W, Chen T. Dynamic modeling and analysis of multi-flexible-link space manipulators under time-varying dynamic boundary conditions[J]. Advances in Space Research, 2024, 74(10): 5224--5243. DOI: \doi{10.1016/j.asr.2024.07.067}.
\bibitem{nenchev1999} Nenchev D N, Yoshida K, Vichitkulsawat P, Uchiyama M. Reaction null-space control of flexible structure mounted manipulator systems[J]. IEEE Transactions on Robotics and Automation, 1999, 15(6): 1011--1023. DOI: \doi{10.1109/70.817666}.
\bibitem{patel2024} Patel D, Damaren C J. Tip and vibration control of space robots using estimated flexible coordinates[J]. Frontiers in Space Technologies, 2024, 5: 1447545. DOI: \doi{10.3389/frspt.2024.1447545}.
\bibitem{zhang2025} Zhang L, Wu J, Yu X, Yu J. Adaptive fuzzy robust $H_\infty$ control for motion and vibration suppression of fully flexible-link space robots[J]. Journal of Computational Methods in Sciences and Engineering, 2025, 25(4): 3019--3033. DOI: \doi{10.1177/14727978251317312}.
\bibitem{ni2019} Ni Z, Liu J, Wu Z, Shen X. Identification of the state-space model and payload mass parameter of a flexible space manipulator using a recursive subspace tracking method[J]. Chinese Journal of Aeronautics, 2019, 32(2): 513--530. DOI: \doi{10.1016/j.cja.2018.05.005}.
\bibitem{forssell1999} Forssell U, Ljung L. Closed-loop identification revisited[J]. Automatica, 1999, 35(7): 1215--1241. DOI: \doi{10.1016/S0005-1098(99)00022-9}.
\bibitem{brunke2022} Brunke L, Greeff M, Hall A W, Yuan Z, Zhou S, Panerati J, Schoellig A P. Safe learning in robotics: From learning-based control to safe reinforcement learning[J]. Annual Review of Control, Robotics, and Autonomous Systems, 2022, 5: 411--444. DOI: \doi{10.1146/annurev-control-042920-020211}.
\bibitem{shi2024} 史玲玲, 肖晓龙, 张晓峰, 范立佳, 单明贺, 田强. 空间机器人在轨组装多模块单元的对接力控制与地面实验[J]. 力学学报, 2024, 56(3): 800--816. DOI: \doi{10.6052/0459-1879-23-458}.
\bibitem{qi2025} 祁伊凡, 单明贺, 田强. 空间多柔体系统缩比等效动力学建模方法研究[J]. 中国科学: 物理学 力学 天文学, 2025, 55(2): 224518. DOI: \doi{10.1360/SSPMA-2024-0302}.
\end{thebibliography}
\endgroup

\subsection{项目的研究内容、研究目标及拟解决的关键科学问题}

\subsubsection{研究内容}

\topic{（1）柔性支撑机器人耦合动力学表征与闭环参数辨识}

研究机械臂运动经安装界面激励支撑结构、结构变形经运动链传递至末端的双向作用关系，建立保留主要惯性耦合项的低阶动力学模型。围绕支撑刚度、阻尼和所携载荷变化，分析不同参数对结构模态、末端误差及控制输入的敏感性，确定需要在线修正的参数组合和适宜保留的模态。结合关节状态、支撑振动及末端位姿测量，研究受限激励下的闭环辨识方法，以物理一致性约束和独立响应校核提高辨识可信度，形成适于控制器调用的动力学模型与状态估计方法。

\topic{（2）参数不确定条件下运动与振动的自适应协同控制}

研究跟踪误差修正与结构振动反馈之间的作用关系，建立考虑支撑响应、驱动能力和任务精度要求的控制方法。以名义动力学前馈和稳定反馈为基础，利用辨识结果修正与主导模态有关的控制参数，通过有界残差补偿减轻可重复模型误差的影响。分析估计误差、参数更新和模态截断对闭环性能的影响，确定控制增益与模型更新速率的适用条件。利用现有实验基础增配柔性支撑试验单元，开展多种支撑和载荷条件下的轨迹运动与停驻实验，验证运动精度和振动衰减的共同改善。

两项内容以可观测动力学变量衔接，前者提供模型及误差范围，后者验证模型修正对控制的作用，地面实验贯穿模型校核与控制验证。

\subsubsection{研究目标与考核指标}

项目总体目标是：揭示柔性支撑机器人作业频段内的主要动力学耦合关系，建立少量敏感参数的闭环辨识方法，形成参数变化条件下兼顾末端运动精度和结构振动的控制方法，并以地面实验验证其有效性和适用范围。

拟采用表\ref{tab:targets}所列指标评价模型有效性、运动误差和残余振动，形成具有可复现测试条件的算法与实验结果。

\begin{table}[htbp]
\centering
\caption{拟定技术指标与评价对象}
\label{tab:targets}
\small
\begin{tabularx}{\linewidth}{P{2.65cm}YP{3.7cm}}
\toprule
评价内容 & 拟达到的指标 & 评价对象与条件 \\
\midrule
动力学模型有效性 & 在预先确定的作业频段内，前两阶可观测主导模态频率相对误差不大于5\%。 & 与独立模态测试结果比较；覆盖3种支撑状态、2种载荷条件。 \\
模型响应预测 & 结构动态位移独立验证数据的归一化均方根预测误差不大于15\%。 & 使用未参与辨识的运动轨迹，统一测点及频段。 \\
运动精度 & 在主要敏感参数存在$\pm10\%$初始建模偏差时，末端位置跟踪均方根误差较固定模型基线降低不小于15\%。 & 比较相同轨迹、执行时间及驱动约束下的重复实验均值。 \\
结构振动 & 运动结束后固定观察窗内，支撑测点动态位移均方根值较固定模型基线降低不小于25\%。 & 与运动精度指标在同组实验中评价，并记录姿态误差和控制输入。 \\
地面验证 & 完成不少于6组支撑与载荷组合，每组开展不少于10次配对重复实验。 & 给出逐工况结果、离散程度及方法适用范围。 \\
\bottomrule
\end{tabularx}
\end{table}

以上指标用于地面原理验证。基线采用经充分整定的名义动力学前馈与稳定反馈控制器，并设置固定参数抑振控制作为补充对照。各控制方法使用相同测量条件；基线参数在标称工况整定后固定。参数偏差取可独立标定的支撑刚度和载荷惯性等敏感参数，分别施加正、负偏差及其代表性组合。指标按各预定工况的重复实验均值分别考核，不以合并平均值掩盖不利工况。

\subsubsection{拟解决的关键科学问题}

\topic{（1）有限观测条件下，柔性支撑响应与机器人运动误差的耦合关系如何被有效辨识}

支撑参数和载荷变化既改变结构振动，也改变机器人运动引起的反作用载荷。不同误差来源可能产生相近的末端响应，使参数修正存在混淆。本项目拟研究参数敏感性、模态可观测性与闭环激励之间的关系，确定能够由有限测量可靠辨识的参数组合，揭示低阶模型保持主要误差传递关系所需的条件。

\topic{（2）模型持续修正时，运动误差校正与结构振动耗散如何协调}

跟踪控制和振动控制通过同一驱动系统作用于刚柔耦合对象，模型更新又改变补偿量和闭环动态。拟研究参数估计误差、更新速率与振动反馈之间的相互影响，确定任务精度与振动约束相容的控制条件，建立在有界模型误差下保持闭环稳定、改善运动与振动性能的方法。

\subsection{拟采取的研究方案及可行性分析}

\subsubsection{总体研究思路与技术路线}

以柔性支撑机器人携载运动为对象建立参考模型，依据参数敏感性和观测条件构造低阶模型，通过闭环数据估计主导参数与柔性状态；将模型修正用于运动与振动控制，分析稳定性及约束条件，并采用独立工况与消融对比检验方法效果。技术路线如图\ref{fig:route}所示。

\begin{figure}[htbp]
\centering
\begin{tikzpicture}[
  box/.style={draw=black!70,rounded corners=2pt,align=center,inner sep=7pt,font=\small,text width=5.8cm,minimum height=1.1cm},
  wide/.style={box,text width=11.6cm},
  arr/.style={-{Latex[length=2mm]},thick,draw=black!75},
  node distance=0.65cm and 0.7cm]
 \node[wide] (object) at (0,0) {\textbf{研究对象与工况}\quad 柔性支撑结构、机械臂及载荷\\既定连接构型\quad 给定运动轨迹\quad 支撑与载荷参数变化};
 \node[box] (model) at (-3.25,-1.9) {\textbf{耦合动力学分析}\\主导模态与误差传递\quad 参数敏感性};
 \node[box] (data) at (3.25,-1.9) {\textbf{闭环数据获取}\\关节状态\quad 结构振动\quad 末端位姿};
 \node[wide] (ident) at (0,-3.8) {\textbf{闭环参数辨识与柔性状态估计}\\可辨识参数组合\quad 物理约束\quad 独立响应校核};
 \node[wide,below=of ident] (control) {\textbf{运动与振动自适应协同控制}\\运动反馈\quad 振动反馈\quad 有界模型补偿\quad 稳定性分析};
 \node[wide,below=of control] (test) {\textbf{多工况地面实验与评价}\\模型预测误差\quad 末端跟踪误差\quad 残余振动\quad 驱动约束};
 \draw[arr] (object.south) -- ++(0,-0.3) -| (model.north);
 \draw[arr] (object.south) -- ++(0,-0.3) -| (data.north);
 \draw[arr] (model.south) -- ++(0,-0.25) -| ($(ident.north)+(-2.3,0)$);
 \draw[arr] (data.south) -- ++(0,-0.25) -| ($(ident.north)+(2.3,0)$);
 \draw[arr] (ident) -- (control);
 \draw[arr] (control) -- (test);
 \draw[arr] (test.east) -- ++(0.6,0) |- (ident.east);
 \node[rotate=90,font=\footnotesize,fill=white,inner sep=2pt] at ($(control.east)+(0.6,0)$) {实验反馈与模型修正};
\end{tikzpicture}
\caption{项目技术路线}
\label{fig:route}
\end{figure}

\subsubsection{耦合动力学建模与闭环辨识方法}

\topic{（1）建立保留界面载荷传递关系的动力学模型}

采用柔性多体系统方法描述机械臂整体运动和支撑的小幅弹性变形，利用有限元或模态模型获得结构参考响应。以整体刚体运动及关节变量$\bm q$、柔性坐标$\bm\eta$组成广义坐标$\bm\zeta$，在给定支撑状态$s$下，将控制设计模型表示为
\begin{equation}
 \bm M_s(\bm\zeta,\bm\theta)\ddot{\bm\zeta}
 +\bm h_s(\bm\zeta,\dot{\bm\zeta},\bm\theta)
 +\begin{bmatrix}\bm 0\\\bm D_s\dot{\bm\eta}+\bm K_s\bm\eta\end{bmatrix}
 =\bm B\bm u+\bm E_s\bm d+\bm r .
 \label{eq:model}
\end{equation}
其中，$\bm M_s$保留运动与弹性变形之间的惯性耦合；$\bm h_s$包括科氏、离心等项；$\bm K_s$和$\bm D_s$为柔性部分的等效刚度与阻尼矩阵；$\bm B$描述实际驱动配置，$\bm d$为外部激励，$\bm r$表示截断和未建模残差。空间模型保留整体刚体运动，地面模型则按实验支承施加相应约束，并计入重力和补偿装置作用。

通过界面反作用载荷和振动能量分析，研究关节加减速、支撑变形与末端误差的传递关系。以作业频段和末端误差敏感性确定保留模态，比较低阶模型与参考模型的频率及运动响应，明确模型适用的变形和载荷范围，以少量参数调整表征结构状态变化。

\topic{（2）构造适于受控运动的参数辨识与状态估计方法}

首先利用结构模态试验和已知载荷运动试验确定参数初值，并进行量纲归一化和灵敏度分析。对高度相关的参数采用组合参数表示，通过观测矩阵和输入输出灵敏度的秩及条件数检验可辨识性；对当前工况下难以区分的参数保持标定值，仅更新能被有效观测的部分。依据主导模态选择加速度或应变测点，并以独立位移测量校核柔性状态估计。

在稳定基线控制下，利用正常轨迹的加减速段和有限幅值的辅助激励获取辨识数据。采用包含噪声模型的闭环预测误差方法进行参数递推，以正刚度、非负阻尼及合理惯性范围约束更新；结合残差相关性和不同数据窗结果，检查噪声与输入相关造成的偏差。以预设的信息量阈值控制更新条件，信息不足时保持参数而不继续漂移。辅助激励在允许的振动和驱动范围内设置，其作用频段覆盖所保留模态附近的有效响应区间。

对重复运动中仍存在的系统性残差，采用低维基函数构造补偿项，并限制其幅值、带宽及变化率。辨识参数和残差补偿分步校准，防止同一误差被重复补偿。所得模型需在未参与辨识的轨迹上通过响应预测校核，方用于后续控制参数修正。

\subsubsection{运动与振动自适应协同控制方法}

\topic{（1）形成兼顾任务误差和振动响应的控制作用}

围绕给定末端轨迹，以名义模型前馈与稳定运动反馈构成基础控制，利用支撑模态速度或其估计量构造振动反馈，并引入有界残差补偿。控制输入采用以下结构组织：
\begin{equation}
 \bm u=\bm u_{\mathrm{tr}}+\bm u_{\mathrm{vib}}+\bm u_{\mathrm{res}},
 \label{eq:control}
\end{equation}
其中，$\bm u_{\mathrm{tr}}$承担运动跟踪，$\bm u_{\mathrm{vib}}$经机械臂与支撑的耦合关系调节振动，$\bm u_{\mathrm{res}}$用于已识别残差的有限补偿。依据辨识结果更新相关模型系数和振动反馈参数，通过输入限幅及变化率约束保持驱动可实现性。

针对欠驱动和任务自由度限制，先分析末端运动与主导模态的可控性，再确定允许跟踪误差、振动反馈强度及控制输入之间的相容范围。控制目标设为给定运动时间内的误差与振动共同降低，避免通过改变任务时长获得表面上的抑振效果。对于不具备独立抑振能力的运动方向，利用受限反馈权重和轨迹加速度分配降低激励，并明确其可达到的性能边界。

\topic{（2）分析模型更新对闭环稳定性与性能的影响}

构造包含运动误差、柔性振动能量和参数估计误差的候选Lyapunov函数，结合观测误差与残差界，研究使闭环误差一致最终有界的充分条件。重点分析参数更新引起的时变项、模态截断及补偿时延，建立更新增益、振动反馈增益与容许误差界之间的关系。在线参数更新采用投影、限速及平滑过渡；当辨识可信度不足时保持已验证参数，并保留稳定基础控制。

控制分析采用包含柔性模态的整体模型，数值验证增加未参与控制设计的高阶模态，检查观测和控制溢出。通过逐项关闭参数更新、振动反馈和残差补偿开展消融实验，区分模型修正与控制结构各自带来的改善，并检验有界残差补偿是否确有必要。

\subsubsection{地面实验方案与评价方法}

依托既有机器人和地面模拟实验条件，配置可更换的柔性梁或桁架支撑单元、标准载荷和同步采集接口。拟选取3种支撑状态与2种载荷条件组成6组工况，开展平滑转运、连续加减速及终点停驻实验。先标定支撑模态、载荷参数和测量噪声，再开展辨识及控制对比。独立校核采用与辨识阶段不同的轨迹参数，防止将训练轨迹上的拟合效果作为模型通用性结论。

支撑结构采用能够产生真实弹性响应的试件，使机器人动作与结构振动形成物理反馈。可重复的力激励或初始变形释放用于设置扰动，不以预设基座位移回放代替抑振验证。末端位姿由外部测量装置独立获取，结构振动以相对静态平衡位置的动态位移评价；针对地面重力和支承附加阻尼，分别建立标定模型并记录其影响。在需要微重力模拟的工况中使用现有平台，结论限定于已验证的动力学特征和参数范围。

对末端位置误差$\bm e_p(k)$、结构测点动态位移$y(k)$及模型预测值$\hat y(k)$，分别采用
\begin{equation}
 E_p=\sqrt{\frac{1}{N}\sum_{k=1}^{N}\|\bm e_p(k)\|^2},\qquad
 E_y=\frac{\sqrt{\sum_k\left[y(k)-\hat y(k)\right]^2}}
 {\sqrt{\sum_k\left[y(k)-\bar y\right]^2}} .
 \label{eq:metrics}
\end{equation}
其中，$\bar y$为验证窗内动态位移均值。运动停止时刻记为$t_f$，以该支撑工况独立标定的一阶固有周期$T_1$确定统一观察窗$T_w=5T_1$，残余振动指标为
\begin{equation}
 V_y=\sqrt{\frac{1}{T_w}\int_{t_f}^{t_f+T_w} y^2(t)\,\mathrm dt} .
 \label{eq:vibration}
\end{equation}
相对改善率按$(E_0-E_1)/E_0$或$(V_0-V_1)/V_0$计算，下标0和1分别代表固定模型基线与本项目方法。验证信号应明显高于测量噪声；接近噪声底的工况报告绝对误差，不据此计算夸大的相对改善率。各方法统一滤波频段、起止时刻和采样配置，同时记录姿态误差、控制输入峰值及重复试验离散程度。

\subsubsection{可行性分析}

柔性多体建模、模态缩减和闭环辨识已有成熟方法，团队的空间操作动力学、结构等效建模及机器人控制积累可支持本项目实施。研究限定于单机器人、有限主导模态和若干既定支撑状态，实验以现有条件的适度改装为主。

针对模态观测不足，依据灵敏度分析调整测点；针对噪声导致的辨识困难，增加受限激励并缩减参数集合；针对在线更新的稳定性影响，采用限速和参数保持；针对高阶模态影响，扩展模型阶次或收缩控制带宽。上述措施可结合阶段试验逐级验证。

\subsection{本项目的特色与创新之处}

\topic{（1）建立面向末端误差传递的柔性支撑动力学辨识方法}

将安装界面的载荷传递、结构主导模态及末端误差联系起来，以可观测性和参数敏感性共同确定低阶模型，使辨识结果直接服务于控制设计。拟阐明受限闭环激励下参数可分辨的条件，形成兼顾物理解释与在线修正能力的耦合动力学表征方法。

\topic{（2）形成辨识误差约束下的运动振动协同控制方法}

将参数估计误差及更新速率纳入运动反馈、振动反馈和有限残差补偿的联合分析，研究有限驱动能力下任务误差与振动耗散的相容条件。拟通过模型更新边界与控制参数的协调设计，在支撑和载荷变化时保持稳定控制，并在同组实验中验证精度与振动指标的共同改善。

\subsection{年度研究计划及预期研究成果}

\subsubsection{年度研究计划}

\topic{2027年1月至6月}
完成典型作业工况和实验参数范围设计，建立柔性支撑机器人参考模型与低阶模型；完成支撑试件设计加工及基础模态测试，分析界面载荷、主导模态与末端误差的关系，确定测点和待辨识参数集合。阶段成果为模型程序、支撑模态标定结果及实验实施方案。

\topic{2027年7月至12月}
完成闭环辨识与柔性状态估计方法研究，校核不同支撑及载荷条件下的模型预测能力；建立运动与振动控制初步方案，完成主要稳定性条件分析和仿真对比。提交年度进展报告，形成论文1篇，组织与重点实验室合作人员的专题研讨并汇报研究进展。

\topic{2028年1月至6月}
完善模型在线修正、振动反馈及残差补偿方法，完成控制程序集成和地面实验调试；开展不同参数偏差下的配对试验与消融对比，修正算法适用范围，形成方法论文和发明专利申请材料。

\topic{2028年7月至12月}
完成6组预定工况的重复实验和独立数据校核，整理逐工况评价结果、程序说明及试验数据；完成论文发表或录用、专利申请和结题报告。与重点实验室开展成果交流及联合评议，提交项目技术总结和实验验证材料。

\subsubsection{预期研究成果与学术交流}

预期形成柔性支撑机器人耦合动力学模型及闭环辨识程序1套，运动与振动自适应协同控制算法1套，配套典型工况数据、评测脚本和使用说明；完成地面原理验证，提交技术总结报告和实验测试报告各1份。围绕上述研究发表或录用学术论文2篇，申请发明专利1项，培养参与研究的研究生2名。

每年与重点实验室研究人员开展至少1次实质性学术交流，结合阶段成果进行模型讨论、实验方案评议及学术报告。拟参加1次相关领域国内学术会议，并利用线上形式与柔性多体动力学及空间机器人领域同行开展交流，国际交流不作为项目完成的前置条件。项目成果的署名、资助标注及知识产权安排按重点实验室开放课题管理要求执行。

\section{研究基础与工作条件}

\subsection{工作基础}

申请人所在团队长期从事大型空间结构动力学与在轨操作研究，在柔性多体系统建模、空间机器人控制和地面实验方面形成了相互支撑的研究积累。与本项目直接相关的基础主要包括以下方面。

\topic{（1）柔性多体动力学与模型缩减基础}

申请人参与了柔性多体系统模态缩减方法研究，相关成果发表于\textit{Multibody System Dynamics}，讨论了模态综合方法及其在整体运动与弹性变形耦合问题中的应用\cite{sonneville2021}。在此基础上，团队开展了空间多柔体系统缩比等效建模研究，结合相似关系和动力学灵敏度分析处理缩比模型设计问题\cite{qi2025}。上述成果可为本项目的主导模态选择、低阶模型校核和实验参数设计提供依据。

\topic{（2）空间机器人动力学与自适应控制基础}

团队已围绕单、双臂空间机器人操作和模块组装开展动力学与控制研究，形成了机器人运动耦合分析、参数不确定性处理和控制实现方面的积累。已发表的多模块单元对接力控制研究建立了相关动力学模型，并通过地面实验检验自适应控制方法\cite{shi2024}。本项目将利用其中的建模、控制实现和实验组织经验，进一步研究柔性支撑条件下的结构响应与运动控制耦合问题。

\topic{（3）空间操作地面实验基础}

团队已建有约$4\,\mathrm m\times10\,\mathrm m$的面向在轨操作的微重力模拟地面实验平台，并具有空间机器人操作相关实验积累。现有平台和已开展的模块组装实验可为本项目提供场地、机器人操作和测试组织基础。本项目拟重点增加柔性支撑试件、同步测量接口及多工况标定环节，形成围绕动力学辨识和振动控制的专项实验条件。

% 不将哈工大指南中的拟完成事项写成已取得成果。
\needinfo{现有控制算法及近期实验结果，包括机器人型号、测试条件和实测指标。}

\subsection{工作条件}

\topic{（1）已具备的条件}

依托单位具有空间结构动力学和机器人控制相关研究团队，已具备柔性多体建模、算法研究和地面实验的基础。现有微重力模拟平台及机器人实验条件能够支撑本项目开展参考模型构建、程序调试和典型操作试验。已有论文成果和实验经验可为测量校核、参数分析及控制性能评价提供方法依据。

\topic{（2）尚需完善的条件及解决途径}

拟根据现有机械臂承载和控制能力，设计可更换的柔性支撑与标准载荷组件，完成结构模态标定；通过共享现有位姿测量和振动测试仪器、适度增配采集与控制接口，建立关节、末端及结构响应的同步记录能力。重点验证控制接口开放程度、采样时序和测量精度，形成可复现的测试流程。试件加工、接口改装和专项标定纳入本项目预算，不新增整套机器人装备。

拟与重点实验室固定研究人员围绕结构模态测试、动力学辨识和闭环控制开展合作，结合双方实际条件安排模型讨论、联合测试及成果交流。具体设备共享范围和试验分工在落实合作人员后确定。

\needinfo{南航重点实验室合作人员姓名、职称及分工；现有控制与测量设备的型号和性能参数。}

\subsection{申请人简介}

\topic{（1）学历与研究工作简历}

单明贺，北京理工大学教授、博士生导师，主要从事空间在轨操作动力学与控制、柔性多体动力学和空间机器人技术研究。2007年9月至2011年7月在哈尔滨工业大学学习，获飞行器制造工程学士学位；2011年9月至2013年7月在哈尔滨工业大学学习，获航空宇航制造工程硕士学位；2013年11月至2018年6月在荷兰代尔夫特理工大学学习，获航天工程博士学位。2018年11月至2020年11月在美国马里兰大学帕克分校从事博士后研究，2021年1月起在北京理工大学工作。

申请人主持国家自然科学基金等科研任务，在空间在轨服务及空间结构动力学相关方向形成了持续研究积累。其研究经历与本项目的柔性多体动力学、机器人运动控制及地面验证紧密相关。

\topic{（2）近期与本项目有关的主要论著}

\begin{enumerate}[label=\arabic*.]
 \begin{samepage}
 \item 祁伊凡, 单明贺, 田强. 空间多柔体系统缩比等效动力学建模方法研究[J]. 中国科学: 物理学 力学 天文学, 2025, 55(2): 224518. DOI: \doi{10.1360/SSPMA-2024-0302}.
 \end{samepage}
 \begin{samepage}
 \item 史玲玲, 肖晓龙, 张晓峰, 范立佳, 单明贺, 田强. 空间机器人在轨组装多模块单元的对接力控制与地面实验[J]. 力学学报, 2024, 56(3): 800--816. DOI: \doi{10.6052/0459-1879-23-458}.
 \end{samepage}
 \begin{samepage}
 \item Sonneville V, Scapolan M, Shan M, Bauchau O A. Modal reduction procedures for flexible multibody dynamics[J]. Multibody System Dynamics, 2021, 51(4): 377--418. DOI: \doi{10.1007/s11044-020-09770-w}.
 \end{samepage}
\end{enumerate}

\topic{（3）学术奖励}

单明贺，国际空间研究委员会青年科学家最佳论文奖（COSPAR Outstanding Paper Award for Young Scientists），2022年。

\needinfo{获奖论文及证书载明信息；中国振动工程学会科学技术二等奖的项目名称、授奖年份、全部受奖人员及本人排序。}

\subsection{项目组主要成员简介}

项目拟按照耦合动力学建模、辨识与控制、实验验证配置人员，由申请人负责总体方案、关键方法研究和项目组织，依托单位研究骨干及研究生承担模型实现、控制程序和实验工作，重点实验室合作人员参与动力学分析、试验方案论证及学术交流。人员名单、个人代表性成果和具体任务须与实际参研安排一致。

\needinfo{按实际参研名单逐人列明简介、承担任务，以及作者和出版信息完整的近期论著、受奖人员和授奖信息完整的学术奖励。}

\subsection{近五年承担科研项目情况}

本项目所依托的既有空间结构动力学和机器人操作工作主要提供理论、算法与实验基础。本项目研究任务为既定连接构型下柔性支撑响应的闭环辨识，以及模型变化条件下的运动振动协同控制，交付对象为对应的模型、算法和专项实验结果。

\needinfo{近五年各科研项目的正式名称、编号、经费来源、起止年月和承担角色，并依据实际任务书逐项说明与本项目的联系、任务边界及成果增量。}

\clearpage
\section{经费预算}

% 预算依据南航2027年度通知所列重点课题20--30万元强度，暂按30万元编制。
% 单价和数量为研究需求测算，尚未取得供应商报价；提交前按实际条件核定。
申请经费30.00万元，围绕柔性支撑试件、测量与控制接口、专项测试及成果产出安排。依托现有机器人和地面实验平台开展研究，优先共享可用仪器设备。经费预算如下。

\begingroup
\small
\begin{longtable}{P{3.15cm}>{\centering\arraybackslash}p{1.35cm}P{9.6cm}}
\toprule
预算科目 & 金额（万元） & 计算根据及理由 \\
\midrule
\endfirsthead
\toprule
预算科目 & 金额（万元） & 计算根据及理由 \\
\midrule
\endhead
\midrule
\multicolumn{3}{r}{续下页}\\
\endfoot
\bottomrule
\endlastfoot
\textbf{合计} & \textbf{30.00} & 服务于两项研究内容及其地面实验验证。 \\
设备费 & 4.00 & 同步采集和实时控制接口改装所需设备及配件1套，按4.00万元测算，用于关节、结构响应与末端测量的时序同步。 \\
材料费 & 6.00 & 3套柔性支撑试件原材料，每套1.20万元，共3.60万元；连接夹具及标准载荷材料1.40万元；传感安装、线缆等耗材1.00万元。 \\
测试化验加工费 & 9.00 & 支撑及夹具加工装配3批，每批1.50万元，共4.50万元；专项模态与测量标定3批，每批0.80万元，共2.40万元；外协动态测量与联合测试3批，每批0.70万元，共2.10万元。共享公用仪器的免费使用不重复计费。 \\
差旅费 & 3.00 & 合作试验、专题交流及参加国内学术会议等约10人次，按每人次0.30万元测算。 \\
会议费 & 0.50 & 项目专题研讨及实验方案评议2次，按每次0.25万元测算。 \\
国际合作与交流费 & 0.00 & 以线上学术交流为主，不安排出国经费。 \\
出版/文献/信息传播/知识产权事务费 & 3.50 & 论文出版相关费用2篇，每篇按1.00万元测算，共2.00万元；发明专利申请及相关事务1.00万元；文献和信息资料0.50万元。 \\
专家咨询费 & 1.00 & 动力学建模及实验方案咨询约5人次，按每人次0.20万元测算；占总经费3.33\%。 \\
劳务费 & 3.00 & 参与建模、程序实现和实验工作的研究生2人，每人累计15个月，每人每月0.10万元；占总经费10.00\%。 \\
\end{longtable}
\endgroup

预算按研究需求测算，设备与测试项目结合现有条件核定，相关支出按实际发生及适用标准执行。专家咨询费和劳务费分别满足申请书模板规定的不超过总经费5\%和10\%的比例要求。

% ==================== 报告正文结束 ====================
\end{document}
